\documentclass[12pt, a4paper]{article}

% ─── Pakete ───────────────────────────────────────────────
\usepackage[utf8]{inputenc}
\usepackage[T1]{fontenc}
\usepackage[ngerman]{babel}
\usepackage{lmodern}
\usepackage[left=2.5cm, right=2.5cm, top=2.5cm, bottom=2.5cm]{geometry}
\usepackage{setspace}
\usepackage{parskip}
\usepackage{enumitem}
\usepackage{booktabs}
\usepackage{tabularx}
\usepackage{array}
\usepackage{hyperref}
\usepackage{amsmath}
\usepackage{csquotes}
\usepackage{xcolor}
\usepackage{titlesec}

\hypersetup{
    colorlinks=true,
    linkcolor=black,
    citecolor=blue!60!black,
    urlcolor=blue!60!black,
}

\onehalfspacing

% ─── Titelseite ───────────────────────────────────────────
\title{%
    \vspace{-1cm}
    \textbf{Exposé zur Masterarbeit}\\[0.8em]
    \Large Effiziente Modellierung multi-stakeholder-basierter\\
    Fragestellungen mittels LLM-Agenten-Simulation:\\[0.3em]
    Entwurf, Implementierung und Evaluation\\
    eines kostenoptimierten Frameworks\\[1em]
    \normalsize\textit{Alternativtitel: A Cost-Efficient Framework for LLM-Based\\
    Multi-Agent Simulations of Stakeholder Negotiations}
}

\author{Janik Häusser}
\date{Mai 2026 — Entwurf}

% ──────────────────────────────────────────────────────────
\begin{document}

\maketitle
\thispagestyle{empty}

\vfill
\begin{center}
    \textit{Entwurf zur Diskussion mit dem Betreuer}
\end{center}
\newpage

\setcounter{page}{1}
\tableofcontents
\newpage

% ══════════════════════════════════════════════════════════
\section{Problemstellung und Motivation}
% ══════════════════════════════════════════════════════════

Large Language Models (LLMs) haben in den letzten Jahren bemerkenswerte Fähigkeiten in der Nachahmung menschlicher Entscheidungsfindung gezeigt. Gleichzeitig wächst das Interesse, diese Modelle als synthetische Akteure in sozialwissenschaftlichen und ökonomischen Experimenten einzusetzen --- etwa zur Simulation von Verhandlungen, Spieltheorie-Experimenten oder Marktdynamiken \cite{horton2023}.

Ein konkretes Beispiel, das die Relevanz solcher Simulationen illustriert, findet sich in der Klimapolitik: Entwässerte Moore auf landwirtschaftlichen Flächen sind eine bedeutende CO\textsubscript{2}-Quelle. Die Wiedervernässung dieser Flächen ist ökologisch wünschenswert, steht jedoch in Konflikt mit den wirtschaftlichen Interessen der Landbesitzer. Verschiedene Stakeholder --- Landwirte mit unterschiedlichen Betriebsgrößen, Risikobereitschaften und ökonomischen Abhängigkeiten, öffentliche Einrichtungen, Umweltorganisationen --- müssen durch geeignete Anreizmechanismen (Subventionen, Kompensationszahlungen, regulatorischer Druck) zu einer Verhaltensänderung bewegt werden. Welche Anreize bei welchen Akteurstypen wirken, ist eine Frage, die sich mittels LLM-basierter Agentensimulation systematisch explorieren lässt --- \textit{bevor} kostenintensive Feldstudien durchgeführt werden.

Bestehende Ansätze für solche Simulationen sind jedoch häufig:
\begin{itemize}
    \item \textbf{ad hoc}: auf ein spezifisches Experiment zugeschnitten, ohne Wiederverwendbarkeit,
    \item \textbf{kostenintensiv}: durch ineffiziente Prompt-Strategien und fehlende Optimierung des API-Verbrauchs,
    \item \textbf{schwer reproduzierbar}: durch mangelnde Standardisierung der Experimental-Pipeline.
\end{itemize}

Es fehlt ein \textbf{generisches, kosteneffizientes Framework}, das es Forschern ermöglicht, solche Multi-Stakeholder-Fragestellungen deklarativ zu definieren und mittels LLM-basierter Agenten systematisch zu simulieren und zu analysieren. Diese Arbeit adressiert diese Lücke und demonstriert das Framework anhand des Anwendungsfalls der Moor-Wiedervernässung.

% ══════════════════════════════════════════════════════════
\section{Zielsetzung}
% ══════════════════════════════════════════════════════════

Ziel dieser Arbeit ist der Entwurf, die Implementierung und die Evaluation eines End-to-End-Frameworks für LLM-basierte Multi-Agenten-Simulationen --- von der deklarativen Szenario-Definition über die Agenten-Modellierung bis zur automatisierten Analyse. Das Framework wird anhand eines konkreten Anwendungsfalls (Moor-Wiedervernässung, siehe Abschnitt~\ref{sec:usecase}) evaluiert und soll folgende Kerneigenschaften aufweisen:

\begin{enumerate}
    \item \textbf{Deklarative Szenario- und Persona-Definition}: Stakeholder, Personas (inkl.\ Hintergrund, Ziele, Constraints), Interaktionsregeln und Anreizmechanismen werden über eine strukturierte Konfiguration (YAML/JSON-Schema) spezifiziert --- ohne Code-Änderungen.

    \item \textbf{Kostenoptimierung}: Systematische Reduktion der API-Kosten durch:
    \begin{itemize}
        \item Prompt-Kompression und Token-Budgetierung,
        \item intelligentes Modell-Routing (günstigere Modelle für einfachere Sub-Tasks),
        \item Prompt-Caching und Response-Memoization,
        \item Structured Output Modes (Reduktion von Parsing-Fehlern und Retries).
    \end{itemize}

    \item \textbf{Performanz und Skalierbarkeit}: Effiziente Parallelisierung, Batching und Rate-Limit-Management für große Kohorten.

    \item \textbf{Wissenschaftliche Rigorosität}: Reproduzierbarkeit durch Seed-Kontrolle, Konfigurationssnapshotting, vollständiges Audit-Trail und standardisierte Analyse-Pipeline.

    \item \textbf{Multi-Provider-Unterstützung}: Abstraktion der LLM-Schnittstelle zur Unterstützung verschiedener Anbieter (OpenAI, Anthropic, Mistral, lokale Modelle via Ollama).
\end{enumerate}

% ══════════════════════════════════════════════════════════
\section{Forschungsfragen}
% ══════════════════════════════════════════════════════════

\subsection*{Hauptfrage}

\begin{quote}
    \textit{Wie lässt sich ein generisches, kosteneffizientes Framework für LLM-basierte Multi-Agenten-Simulationen von Stakeholder-Verhandlungen entwerfen, implementieren und evaluieren?}
\end{quote}

\subsection*{Teilfragen}

\begin{description}[style=nextline, leftmargin=1.5cm]
    \item[F1 (Effizienz)] Welche Prompt-Engineering- und API-Optimierungsstrategien reduzieren die Kosten pro Simulation signifikant, ohne die Qualität der Agenten-Entscheidungen zu beeinträchtigen?

    \item[F2 (Architektur)] Wie muss die Framework-Architektur gestaltet sein, damit verschiedene Stakeholder-Szenarien (Verhandlung, iteratives Spiel, Multi-Stakeholder-Konflikte) durch reine Konfigurationsänderungen abgebildet werden können?

    \item[F3 (Validität)] Produzieren die simulierten Agenten plausible, konsistente Entscheidungen im Kontext bekannter spieltheoretischer Szenarien und eines realen Anwendungsfalls (Moor-Wiedervernässung)? Wie sensitiv sind die Ergebnisse gegenüber Modellwahl, Temperatur und Prompt-Variationen?

    \item[F4 (Skalierung)] Wie skalieren Kosten und Laufzeit mit der Anzahl der Agenten, Runden und Komplexität der Szenarien, und welche architektonischen Entscheidungen beeinflussen dies?
\end{description}

% ══════════════════════════════════════════════════════════
\section{Methodik}
% ══════════════════════════════════════════════════════════

\subsection{Design Science Research}

Die Arbeit folgt dem \textbf{Design Science Research}-Paradigma \cite{hevner2004}: Entwurf eines IT-Artefakts (Framework), Implementierung und systematische Evaluation.

\subsection{Technische Umsetzung}

\subsubsection*{Phase~1 --- Analyse und Anforderungen}
\begin{itemize}
    \item Systematische Analyse bestehender Frameworks (AutoGen, CrewAI, LangGraph, CAMEL) und Abgrenzung.
    \item Anforderungserhebung anhand von Referenz-Szenarien aus der experimentellen Ökonomie.
    \item Domänenanalyse des Anwendungsfalls Moor-Wiedervernässung: Identifikation relevanter Stakeholder-Typen, Anreizmechanismen und Entscheidungsparameter.
    \item Analyse des bestehenden ABM-Prototypen (Gift-Exchange-Simulation mit Mistral).
\end{itemize}

\subsubsection*{Phase~2 --- Architektur und Implementierung}

Entwurf einer modularen Architektur mit folgenden Komponenten:
\begin{description}[style=nextline, leftmargin=1cm]
    \item[Scenario Engine] Parst deklarative Szenario-Definitionen inkl.\ Stakeholder, Anreize und Constraints.
    \item[Persona Manager] Instantiiert heterogene Agenten aus deklarativen Persona-Definitionen (Hintergrund, Ziele, Risikobereitschaft, Entscheidungslogik).
    \item[Interaction Orchestrator] Steuert Kommunikationstopologien und Rundenlogik (1:1-Verhandlung, Gruppeninteraktion, iterative Runden).
    \item[LLM Abstraction Layer] Provider-agnostische API-Schicht mit Retry, Caching, Routing.
    \item[Cost Controller] Token-Tracking, Budget-Enforcement, Modell-Routing-Logik.
    \item[Analysis Pipeline] Standardisierte Metriken, Visualisierungen, Export.
\end{description}

Implementierung in Python mit Fokus auf Erweiterbarkeit und Testbarkeit. Das Framework folgt einem dreischichtigen Aufbau:

\begin{enumerate}
    \item \textbf{Input-Schicht}: Drei YAML-Dateien --- Szenario-Definition (Rollen, Interaktionsregeln, Anreize, Constraints), Persona-Definitionen (Hintergrund, Ziele, Risikoprofil, Entscheidungsfaktoren) und Run-Konfiguration (LLM-Provider, Budget, Seed, Parallelisierung). Alle Dateien werden gegen JSON-Schemata validiert.
    \item \textbf{Core Engine}: Scenario Engine, Persona Manager, Interaction Orchestrator, LLM Gateway (Multi-Provider mit Routing und Caching), Cost Controller und State Manager (immutable per Round).
    \item \textbf{Output-Schicht}: Strukturierte Ergebnisse --- Decision Logs (JSONL), Cost Logs, Raw Responses, Reproduzierbarkeits-Manifest (Config-Hash, Seed, Framework-Version, Modellversionen) sowie aggregierte Analysen.
\end{enumerate}

Eine ausführliche technische Spezifikation mit Schema-Beispielen, Datenflussdiagrammen und Kostenoptimierungsstrategien liegt als separates Architektur-Dokument vor.

\subsubsection*{Phase~3 --- Evaluation}
\begin{itemize}
    \item \textbf{Anwendungsszenario}: Simulation der Moor-Wiedervernässungs-Verhandlung mit verschiedenen Landwirt-Personas und Anreizmechanismen. Qualitative Analyse der Agenten-Entscheidungen auf Plausibilität und Konsistenz.
    \item \textbf{Benchmark-Szenarien}: Prisoner's Dilemma, Ultimatum Game, Public Goods Game --- Vergleich mit theoretischen Gleichgewichten als Validierung der Agenten-Qualität.
    \item \textbf{Kostenvergleich}: Messung der API-Kosten (Tokens, \$) mit vs.\ ohne Optimierungen über identische Szenarien.
    \item \textbf{Sensitivitätsanalyse}: Variation von Modell, Temperatur, Prompt-Strategie und Kohortengrößen.
    \item \textbf{Reproduzierbarkeitstest}: Wiederholung identischer Konfigurationen und statistische Auswertung der Varianz.
\end{itemize}

\subsection{Evaluationskriterien}

\begin{table}[h!]
    \centering
    \begin{tabularx}{\textwidth}{lX}
        \toprule
        \textbf{Kriterium} & \textbf{Metrik} \\
        \midrule
        Kosteneffizienz        & Tokens/Simulation, \$/Simulation, Tokens/Entscheidung \\
        Performanz             & Latenz/Runde, Gesamtdauer, Durchsatz (Simulationen/Stunde) \\
        Validität              & Abweichung von theoretischen Gleichgewichten (MSE, Korrelation) \\
        Reproduzierbarkeit     & Varianz über identische Runs (Konfidenzintervalle) \\
        Generalisierbarkeit    & Anzahl unterstützter Szenariotypen ohne Code-Änderung \\
        Usability              & Lines of Config pro Szenario, Time-to-First-Result \\
        \bottomrule
    \end{tabularx}
    \caption{Evaluationskriterien und zugehörige Metriken}
    \label{tab:eval}
\end{table}

% ══════════════════════════════════════════════════════════
\section{Primäres Anwendungsszenario: Moor-Wiedervernässung}
\label{sec:usecase}
% ══════════════════════════════════════════════════════════

Zur Demonstration und Evaluation des Frameworks dient ein konkreter Anwendungsfall aus dem Bereich Klimapolitik und Landnutzung:

\textbf{Kontext.} In Deutschland und Europa sind große Moorflächen für landwirtschaftliche Nutzung entwässert worden. Entwässerte Moore emittieren erhebliche Mengen CO\textsubscript{2} und sind damit relevante Treibhausgas-Quellen. Die Wiedervernässung (Rewetting) dieser Flächen ist ein effektiver Beitrag zum Klimaschutz, erfordert jedoch, dass Landbesitzer ihre bisherige Bewirtschaftungsform aufgeben oder anpassen.

\textbf{Stakeholder.} Das Szenario umfasst mehrere Akteursgruppen mit heterogenen Interessen:
\begin{itemize}
    \item \textbf{Landwirte} (diverse Personas): konventionelle Großbetriebe mit hoher Abhängigkeit von Ackerbau, kleinere Betriebe mit Diversifizierungspotenzial, ökologisch orientierte Landwirte, risikoscheue vs.\ risikofreudige Betriebsleiter.
    \item \textbf{Öffentliche Einrichtungen}: Kommunen, Landesbehörden oder Bundesministerien, die Kompensationszahlungen, Subventionen oder regulatorische Vorgaben anbieten.
    \item \textbf{Umweltorganisationen}: NGOs, die moralischen und öffentlichen Druck ausüben, aber auch Kooperationsangebote machen.
\end{itemize}

\textbf{Fragestellung.} Welche Kombination aus Anreizmechanismen (Kompensationshöhe, Vertragslaufzeit, Umstiegsförderung, Paludikultur-Beratung, regulatorischer Druck) führt bei welchen Landwirt-Personas zur Akzeptanz der Wiedervernässung? Wie verändern sich die Ergebnisse unter Variation der Anreizparameter?

\textbf{Einordnung.} Das Szenario dient primär als \textit{technische Demonstration} des Frameworks. Die inhaltliche Validität der Ergebnisse (d.\,h.\ ob LLM-Agenten reale Landwirte korrekt abbilden) ist \textit{nicht} der Kern dieser Informatik-Arbeit, wird aber als Plausibilitätsprüfung diskutiert und als Anknüpfungspunkt für domänenspezifische Folgearbeiten benannt.

% ══════════════════════════════════════════════════════════
\section{Abgrenzung zu bestehenden Arbeiten}
% ══════════════════════════════════════════════════════════

\begin{table}[h!]
    \centering
    \small
    \begin{tabularx}{\textwidth}{lp{3.8cm}X}
        \toprule
        \textbf{Framework / Arbeit} & \textbf{Fokus} & \textbf{Abgrenzung dieser Arbeit} \\
        \midrule
        AutoGen \cite{wu2023}
            & Multi-Agent Conversations für Software-Tasks
            & Kein Fokus auf wissenschaftliche Simulation, keine Kostenoptimierung \\
        CrewAI
            & Task-orientierte Agent-Teams
            & Workflow-Orchestrierung, nicht wissenschaftliche Reproduzierbarkeit \\
        LangGraph
            & Zustandsbasierte Agent-Graphen
            & Generisches Tool, kein domänenspezifisches Szenario-Schema \\
        CAMEL \cite{li2023}
            & Kommunikative Agenten
            & Role-Playing fokussiert, keine ökonomische Evaluation \\
        Horton \cite{horton2023}
            & LLMs als ökonomische Agenten
            & Einzelexperimente, kein Framework \\
        ABM-Prototyp (Betreuer)
            & Gift-Exchange-Simulation
            & Domänenspezifisch, Single-Provider, keine Kostenoptimierung \\
        \midrule
        \textbf{Diese Arbeit}
            & \textbf{Generisches, kosten\-optimiertes Simulations-Framework}
            & \textbf{Deklarativ, Multi-Provider, kostenoptimiert, reproduzierbar} \\
        \bottomrule
    \end{tabularx}
    \caption{Abgrenzung zu verwandten Arbeiten}
    \label{tab:related}
\end{table}

% ══════════════════════════════════════════════════════════
\section{Vorläufige Gliederung der Arbeit}
% ══════════════════════════════════════════════════════════

\begin{enumerate}
    \item \textbf{Einleitung}
    \begin{enumerate}[label=\arabic{enumi}.\arabic*]
        \item Motivation und Problemstellung
        \item Zielsetzung und Forschungsfragen
        \item Aufbau der Arbeit
    \end{enumerate}

    \item \textbf{Grundlagen}
    \begin{enumerate}[label=\arabic{enumi}.\arabic*]
        \item Large Language Models: Architektur, APIs, Kostenmodelle
        \item Agentenbasierte Modellierung (ABM)
        \item Spieltheorie und Stakeholder-Verhandlungen
        \item Bewertung von LLM-Agenten: Validität, Reproduzierbarkeit
        \item Domänenkontext: Moorentwässerung, Klimapolitik und Anreizmechanismen
    \end{enumerate}

    \item \textbf{Stand der Forschung}
    \begin{enumerate}[label=\arabic{enumi}.\arabic*]
        \item LLMs als simulierte ökonomische Akteure
        \item Multi-Agenten-Frameworks (AutoGen, CrewAI, LangGraph, CAMEL)
        \item Kostenoptimierung bei LLM-Anwendungen
        \item Forschungslücke und Beitrag dieser Arbeit
    \end{enumerate}

    \item \textbf{Anforderungsanalyse}
    \begin{enumerate}[label=\arabic{enumi}.\arabic*]
        \item Funktionale Anforderungen
        \item Nicht-funktionale Anforderungen (Kosten, Performanz, Reproduzierbarkeit)
        \item Anwendungsszenario Moor-Wiedervernässung: Stakeholder und Anreize
        \item Benchmark-Szenarien aus der Spieltheorie
    \end{enumerate}

    \item \textbf{Framework-Entwurf}
    \begin{enumerate}[label=\arabic{enumi}.\arabic*]
        \item Architekturübersicht
        \item Szenario-Definitionssprache (Schema-Entwurf)
        \item Agent-Lifecycle und Persona-Modellierung
        \item Interaktions-Orchestrierung und Topologien
        \item LLM-Abstraktionsschicht und Multi-Provider-Integration
        \item Kostenoptimierungsstrategien
        \item Analyse- und Reporting-Pipeline
    \end{enumerate}

    \item \textbf{Implementierung}
    \begin{enumerate}[label=\arabic{enumi}.\arabic*]
        \item Technologie-Stack und Projektstruktur
        \item Kernkomponenten und Schnittstellen
        \item Prompt-Engineering-Strategien
        \item Caching, Batching und Parallelisierung
        \item Testabdeckung und Qualitätssicherung
    \end{enumerate}

    \item \textbf{Evaluation}
    \begin{enumerate}[label=\arabic{enumi}.\arabic*]
        \item Experimentelles Setup
        \item Anwendungsszenario: Moor-Wiedervernässungs-Simulation
        \item Benchmark-Ergebnisse: Validität in spieltheoretischen Szenarien
        \item Kostenanalyse: Optimiert vs.\ Baseline
        \item Performanz-Benchmarks: Latenz und Durchsatz
        \item Reproduzierbarkeitsanalyse
        \item Diskussion und Limitationen
    \end{enumerate}

    \item \textbf{Fazit und Ausblick}
    \begin{enumerate}[label=\arabic{enumi}.\arabic*]
        \item Zusammenfassung der Ergebnisse
        \item Beantwortung der Forschungsfragen
        \item Limitationen
        \item Ausblick und zukünftige Arbeiten
    \end{enumerate}
\end{enumerate}

% ══════════════════════════════════════════════════════════
\section{Vorläufiger Zeitplan}
% ══════════════════════════════════════════════════════════

\begin{table}[h!]
    \centering
    \begin{tabularx}{\textwidth}{llX}
        \toprule
        \textbf{Phase} & \textbf{Zeitraum} & \textbf{Aufgaben} \\
        \midrule
        1.\ Analyse
            & Monat 1--2
            & Literaturrecherche, Anforderungsanalyse, Domänenanalyse Moor-Szenario, Benchmark-Definition, Analyse des Prototypen \\
        2.\ Entwurf
            & Monat 2--3
            & Architektur-Design, Schema-Entwurf, Technologieauswahl \\
        3.\ Implementierung
            & Monat 3--5
            & Kernframework, Persona-Modellierung, LLM-Integration, Kostenoptimierungen, Tests \\
        4.\ Evaluation
            & Monat 4--5
            & Moor-Szenario-Simulation, Benchmark-Experimente, Kostenvergleiche, Sensitivitätsanalysen \\
        5.\ Verschriftlichung
            & Monat 5--6
            & Schreiben, Review, Überarbeitung \\
        \bottomrule
    \end{tabularx}
    \caption{Vorläufiger Zeitplan (Annahme: 6 Monate Bearbeitungszeit)}
    \label{tab:zeitplan}
\end{table}

% ══════════════════════════════════════════════════════════
\section{Erwarteter Beitrag}
% ══════════════════════════════════════════════════════════

\begin{enumerate}
    \item \textbf{Praxisbeitrag}: Ein Open-Source-Framework, das Forschern ermöglicht, Multi-Stakeholder-Simulationen über deklarative Konfiguration durchzuführen --- demonstriert am Anwendungsfall der Moor-Wiedervernässung.
    \item \textbf{Wissenschaftlicher Beitrag}: Systematische Evaluation von Kostenoptimierungsstrategien für LLM-Agentensysteme mit quantitativen Benchmarks.
    \item \textbf{Methodischer Beitrag}: Entwurfsmuster und Best Practices für reproduzierbare LLM-basierte Experimentalsimulationen, einschließlich eines deklarativen Persona-Schemas.
    \item \textbf{Anwendungsbeitrag}: Erste explorative Ergebnisse zur Frage, welche Anreizmechanismen bei verschiedenen Landwirt-Typen zur Akzeptanz von Wiedervernässungsmaßnahmen führen können.
\end{enumerate}

% ══════════════════════════════════════════════════════════
\section{Benötigte Ressourcen}
% ══════════════════════════════════════════════════════════

\begin{itemize}
    \item LLM-API-Zugang (OpenAI, Anthropic, Mistral) --- geschätztes Budget: 50--200\,€ für Evaluation.
    \item Compute: Lokaler Rechner für Framework-Entwicklung; ggf.\ GPU für lokale Modelle (Ollama).
    \item Zugang zum bestehenden Prototyp-Repository des Betreuers.
    \item Zugang zu relevanter Literatur (Universitätsbibliothek).
\end{itemize}

% ══════════════════════════════════════════════════════════
% Literaturverzeichnis
% ══════════════════════════════════════════════════════════

\begin{thebibliography}{99}

\bibitem{horton2023}
Horton, J.\,J. (2023).
\textit{Large Language Models as Simulated Economic Agents: What Can We Learn from Homo Silicus?}
NBER Working Paper No.\ 31122.

\bibitem{park2023}
Park, J.\,S., O'Brien, J.\,C., Cai, C.\,J., Morris, M.\,R., Liang, P. \& Bernstein, M.\,S. (2023).
Generative Agents: Interactive Simulacra of Human Behavior.
In: \textit{Proceedings of UIST~'23}.

\bibitem{aher2023}
Aher, G.\,V., Arriaga, R.\,I. \& Kalai, A.\,T. (2023).
Using Large Language Models to Simulate Multiple Humans and Replicate Human Subject Studies.
In: \textit{Proceedings of ICML~'23}.

\bibitem{li2023}
Li, G., Hammoud, H.\,A.\,A.\,K., Itani, H., Khizbullin, D. \& Ghanem, B. (2023).
CAMEL: Communicative Agents for ``Mind'' Exploration of Large Language Model Society.
In: \textit{Proceedings of NeurIPS~'23}.

\bibitem{wu2023}
Wu, Q., Bansal, G., Zhang, J. et al. (2023).
AutoGen: Enabling Next-Gen LLM Applications via Multi-Agent Conversation.
\textit{arXiv preprint} arXiv:2308.08155.

\bibitem{hevner2004}
Hevner, A.\,R., March, S.\,T., Park, J. \& Ram, S. (2004).
Design Science in Information Systems Research.
\textit{MIS Quarterly}, 28(1), 75--105.

\bibitem{teubner2025}
Teubner, T. \& Camacho, S. (2025).
[Referenz für Gift-Exchange-Studien mit LLMs --- Quelle beim Betreuer erfragen].

\bibitem{chen2023}
Chen, L., Zaharia, M. \& Zou, J. (2023).
FrugalGPT: How to Use Large Language Models While Reducing Cost and Improving Performance.
\textit{arXiv preprint} arXiv:2305.05176.

\end{thebibliography}

\end{document}
